% Generated by tools/textbook/export_listings.js from tools/concepts/assessments.json: do not edit.
\begin{chapstatus}{Il filo del capitolo}
Il ra salvato nella Lezione 02 è ora spiegato: lo stack e la convenzione di chiamata sono regole che permettono a funzioni scritte da persone diverse di chiamarsi. Prossimo passo: come un'istruzione è codificata ed eseguita, e come il processore viene interrotto.
\par\medskip\textbf{Biglietto d'uscita} (da due a quattro minuti, da soli, su carta):
\begin{enumerate}[nosep,leftmargin=2em]
\item Seguite add t2, t0, t1 nello schema a blocchi di un processore RV32: nominate, in ordine, i blocchi che usa e che cosa fa ciascuno. Quale blocco non usa?
\item Una funzione comincia con addi sp, sp, -16 e sw ra, 12(sp). Dove si trova ora ra, e che cosa deve fare la funzione prima di ret?
\end{enumerate}
\end{chapstatus}
